\documentclass{article}
\usepackage[UTF8]{ctex}
\usepackage{amsmath,amssymb,amsfonts}
\usepackage{enumitem}
\usepackage{geometry}
\geometry{a4paper, margin=1in}
% === 错题本排版优化（不参与题目标识） ===
\setlength{\parindent}{0pt}
\setlength{\parskip}{0.35em plus 0.10em minus 0.08em}
\setlength{\emergencystretch}{3em}
\setlist{itemsep=0.25em, topsep=0.40em, parsep=0pt}
\allowdisplaybreaks[2]
% === 错题本排版优化结束 ===
\begin{document}
\section*{陕西师范大学2026年高等代数真题}
\section*{179．陕西师范大学2026年高等代数试卷}

\subsection{题目1}
1．（10 分）设 \(f(x)=x^{4}+x^{3}+3 x^{2}-4 x-1, g(x)=x^{3}+x^{2}-x-1\) ，求（ \(f(x), g(x)\) ），并求 \(u(x), v(x)\) ，使得 \(u(x) f(x)+v(x) g(x)=(f(x), g(x))\) ．

\subsection{题目2}
2．（10 分）设 4 级矩阵

\[
A=\left(\begin{array}{c}
\alpha \\
2 \gamma_{2} \\
3 \gamma_{3} \\
4 \gamma_{4}
\end{array}\right), B=\left(\begin{array}{c}
\beta \\
\gamma_{2} \\
\gamma_{3} \\
\gamma_{4}
\end{array}\right)
\]

其中 \(\alpha, \beta, \gamma_{2}, \gamma_{3}, \gamma_{4}\) 均为4维行向量，且已知 \(|A|=2800,|B|=100\) ，试计算行列式 \(|A-B|\) 。

\subsection{题目3}
3．（10 分）在齐次线性方程组

\[
\left\{\begin{array}{c}
a_{11} x_{1}+a_{12} x_{2}+\cdots+a_{1 n} x_{n}=0, \\
a_{21} x_{1}+a_{22} x_{2}+\cdots+a_{2 n} x_{n}=0, \\
\cdots \cdots \\
a_{n-1,1} x_{1}+a_{n-1,2} x_{2}+\cdots+a_{n-1, n} x_{n}=0
\end{array}\right.
\]

中，证明

\[
\begin{gathered}
x_{1}=\left|\begin{array}{cccc}
a_{12} & a_{13} & \cdots & a_{1 n} \\
a_{22} & a_{23} & \cdots & a_{2 n} \\
\vdots & \vdots & & \vdots \\
a_{n-1,2} & a_{n-1,3} & \cdots & a_{n-1, n}
\end{array}\right|, x_{2}=-\left|\begin{array}{cccc}
a_{11} & a_{13} & \cdots & a_{1 n} \\
a_{21} & a_{23} & \cdots & a_{2 n} \\
\vdots & \vdots & & \vdots \\
a_{n-1,1} & a_{n-1,3} & \cdots & a_{n-1, n}
\end{array}\right|, \\
\cdots, x_{n}=(-1)^{n-1}\left|\begin{array}{cccc}
a_{11} & a_{12} & \cdots & a_{1, n-1} \\
a_{21} & a_{22} & \cdots & a_{2, n-1} \\
\vdots & \vdots & & \vdots \\
a_{n-1,1} & a_{n-1,2} & \cdots & a_{n-1, n-1}
\end{array}\right|
\end{gathered}
\]

是方程组的解，且若这个解不为零，则方程组的任意解可由它乘以某数得到。

\subsection{题目4}
4．（10 分）设 \(A \in M_{m \times n}(P)\) ，证明：\(A\) 的秩等于 \(r\) 的充要条件是存在秩为 \(r\) 的 \(m \times r\) 矩阵 \(M\) 和\(r \times n\) 矩阵 \(N\) ，使得 \(A=M N\) 。

\subsection{题目5}
5．（10 分）若 \(A=\left(a_{i j}\right)_{n \times n}\) 是正定阵，证明：

\[
0<|A| \leq a_{11} a_{22} \cdots a_{n n}
\]

其中 \(a_{i i}\) 为 \(A\) 的主对角线上的元素，\(i=1,2, \cdots, n\) 。

\subsection{题目6}
6．（20 分）设 \(A\) 是数域 \(P\) 上的 \(n\) 级幂等矩阵，即 \(A^{2}=A\) ．齐次线性方程组 \(A X=0\) 的解空间为 \(W_{1}\) ，\((A-E) X=0\) 的解空间为 \(W_{2}\) ，证明：\(P^{n}=W_{1} \oplus W_{2}\) 。

\subsection{题目7}
7．（20 分）设 \(V\) 是数域 \(P\) 上的 3 维线性空间，线性变换 \(f: V \rightarrow V\) 在 \(V\) 的基 \(e_{1}, e_{2}, e_{3}\) 下的矩阵为

\[
A=\left(\begin{array}{lll}
4 & 6 & -15 \\
1 & 3 & -5 \\
1 & 2 & -4
\end{array}\right)
\]

问 \(f\) 可否在 \(V\) 的某个基下的矩阵为 \(B=\left(\begin{array}{ccc}1 & -3 & 3 \\ -2 & -6 & 13 \\ -1 & -4 & 8\end{array}\right)\) ，为什么？

\subsection{题目8}
8．（20 分）设 \(A\) 是实数域 \(\mathbb{R}\) 上的 \(n\) 级可逆阵，证明：存在实数域上的 \(n\) 级正定阵 \(P\) 和 \(n\) 级正交阵 \(Q\)，使得 \(A=P Q\) ，并且这一分解式是惟一的．

\subsection{题目9}
9．（20 分）设 \(A=\left(\begin{array}{ccc}1 & 0 & 2 \\ 0 & -1 & 1 \\ 0 & 1 & 0\end{array}\right)\) ，求 \(2 A^{8}-3 A^{5}+A^{4}+A^{2}-4 E\) ．

\subsection{题目10}
10．（20分）设 \(n\) 级实矩阵 \(A\) 与 \(B\) 都可对角化，即存在可逆矩阵 \(S\) 与 \(T\) ，使得 \(S^{-1} A S\) 与 \(T^{-1} B T\)都是对角矩阵．证明：当 \(A B=B A\) 时，存在可逆矩阵 \(C\) ，使得 \(A\) 与 \(B\) 可同时对角化，即 \(C^{-1} A C\)与 \(C^{-1} B C\) 都是对角阵。
\end{document}
